%==============================================================
% ABSTRACT + INDEX TERMS -- journal version.
%
% Every figure quoted here is traceable:
%   100 percent metadata-only            artefact_control.csv
%   1.03 point cost of re-encoding       95.38 - 94.35, same file
%   1,000 of 1,000 byte-identical        split_ceilings.csv
%   AUROC 0.4832 / 0.5165 / 0.5094       textsplit_proof.csv
%   95.38 / 94.35, text-only below 78    ablation_grouped_summary.csv + above
%   94.35 -> 84.9 under shift (9.45)     big_sealed_results_v2.csv
%   n = 1,308 sealed                     big_sealed_results_v2.csv
%   +0.0240 / +0.0273 AUROC, 8 seeds     adaptation_final.csv
%==============================================================

\begin{abstract}
Fake news detection for low-resource languages such as Bangla depends on
synthetically constructed benchmarks, and a constructed benchmark is assembled
by a pipeline that leaves traces. This paper reports an image-level detector
for manipulated Bangladeshi news imagery together with the controls needed to
say what its numbers mean. Two properties of the BanFakeMM benchmark change the
interpretation of any result obtained on it. First, file metadata
alone---resolution, format and colour mode, with no pixel content---separates
manipulated from authentic images at 100 percent accuracy; re-encoding every
image through one identical pipeline removes the artefact at a cost of
1.03 accuracy points, which establishes that the remaining signal lies in the
content; a weaker version of the same shortcut (77 to 83 percent) appears in a
second Bangla corpus. Second, the benchmark's five attack types do not pose a single task:
in three of them every one of the 1,000 manipulated images is byte-identical to
an authentic image, so an image-only model is undecidable there, and its
measured AUROC on those splits is 0.48 to 0.52. We scope the evaluation to the
two splits in which the manipulation reaches the pixels. Within that scope, a
per-modality ablation under duplicate-aware grouped cross-validation shows CLIP
visual features dominating at 95.38 percent---94.35 with the artefact
removed---while every text-only configuration saturates below 78 percent; OCR
and the caption-fusion pathway are removed on evidence rather than by
assumption. Transfer to an unseen manipulation style costs roughly nine points
of accuracy on a 1,308-image sealed partition, and supervised adaptation on
300 target-style images improves ranking by 0.025 AUROC, with confidence
intervals excluding zero across eight seeds and both encodings. Two conclusions
we previously published about this system---that the shift penalty was largely
a calibration effect, and that supervised adaptation added nothing---are
corrected here, together with the measurements that force the correction.
\end{abstract}

\begin{IEEEkeywords}
Bangla fake news detection, multimodal misinformation, CLIP, distribution
shift, dataset artefacts, shortcut learning, near-duplicate leakage,
evaluation methodology, low-resource language processing
\end{IEEEkeywords}
